\section{Conclusiones}
\label{sec:conclusiones}
\textit{Responsable: Mariano, con aportes de todo el equipo.}

\subsection{Lo que se cumplió}
\begin{itemize}
  \item Un LMS multi-tenant para tres instituciones con seis microservicios serverless, cada uno con su
        API, sus Lambdas y su tabla DynamoDB.
  \item Comunicación por eventos con outbox transaccional, EventBridge Pipes, una cola con DLQ por
        consumidor y consumidores idempotentes; orden por curso para las notas con SNS/SQS FIFO.
  \item Autenticación con Cognito y aislamiento por tenant tomado siempre del token.
  \item Despliegue reproducible con Serverless Framework en cualquier cuenta del Learner Lab, verificado
        con los seis microservicios en staging y las alarmas en verde.
\end{itemize}

\subsection{Lecciones aprendidas}
\begin{itemize}
  \item Probar pronto en el Learner Lab lo que no se sabe si está permitido: EventBridge Pipes sí
        funcionó y simplificó el diseño; Schema Registry y el visor de X-Ray no, y se buscaron
        alternativas.
  \item Las pruebas de carga encuentran problemas de diseño de datos que las pruebas funcionales no
        ven (el listado de cursos sin índice).
\end{itemize}

\pendiente{Que cada integrante aporte una lección de su parte.}

\subsection{Trabajo futuro}
Propagar \texttt{traceparent} en los eventos para encadenar la traza entre microservicios, índices
secundarios para los listados por usuario y las fases 2 a 4 del \textit{roadmap} (quizzes, foros,
analítica e integraciones externas).
